\documentclass[11pt]{article}
\usepackage[margin=1in]{geometry}
\usepackage{xcolor}
\usepackage{enumitem}
\usepackage{amsmath,amssymb,mathtools}
\usepackage{hyperref}

\definecolor{revcolor}{rgb}{0,0,0.7}
\newcommand{\RC}[1]{\textcolor{revcolor}{\textit{#1}}}
\newcommand{\resp}[1]{\smallskip\noindent #1\medskip}

\title{Response to Reviewer Comments\\[4pt]
\large ``Effective Inversely Proportional Hypermutations\\for Unimodal and Multimodal Optimisation''}
\author{}
\date{}

\begin{document}
\maketitle

We thank the reviewer for a careful and constructive reading of the manuscript.
Below we address each comment point by point. All changes in the revised
manuscript are typeset in {\textcolor{blue}{blue}}.

%% ======================================================================
\section*{Proof Tightening: Theorems 2, 4, 5, 6 \& 7}
%% ======================================================================

\subsection*{Theorem 2 (OneMax, expoF/expoHD)}

\RC{The upper bound can be reduced by a logarithmic factor by splitting the sum
into three ranges $[1,\, n/\!\ln n]$, $[n/\!\ln n,\, \delta n]$, and
$[\delta n,\, n/2 + n^{2/3}]$. For the lower bound, with probability at least
$1/2$ we start with $i \ge n/2$, and the sum over $[n/4,\, n/2]$ can be bounded
via a geometric progression, giving $\sqrt{n}/\!\log n$.}

\resp{We have followed the reviewer's suggestion in full.
The proof of Theorem~2 has been completely rewritten.
The upper bound now splits the sum into the three indicated ranges, yielding
$O(n^{3/2}/\!\log n)$. The lower bound sums over $[n/4,\, n/2]$ using the
geometric-series evaluation as suggested, giving $\Omega(n^{3/2}/\!\log n)$.
The theorem statement has been updated to $\Theta(n^{3/2}/\!\log n)$.}

\subsection*{Theorem 4 (LeadingOnes, expoF)}

\RC{Why do you restrict yourself to the interval $[n/2+\varepsilon n/2,\;
n/2+\varepsilon n]$? With probability $1/2$ we start with $i=0$ and hence we
can sum over all $i\in[0..n-1]$. The sum $\sum n\,i\,n^{-i/n}$ can be estimated
via $\sum i\,x^i$ with $x=n^{-1/n}<1$, which gives $n^3/\!\log n$.}

\resp{We have adopted the reviewer's approach: the lower bound now sums
$E(f_i)$ over the full range $i\in[0,\,n-1]$ and evaluates the resulting series
via the identity $\sum i\,x^i = x/(1-x)^2$ with $x = n^{-1/n}$.
The tightened result is $\Theta(n^3/\!\log^2 n)$ (matching upper and lower
bounds). The theorem statement and proof have been updated accordingly.}

\subsection*{Theorem 5 (LeadingOnes, expoHD)}

\RC{(Implicit from the tightening programme for Theorems~2 and~4.)}

\resp{Following the same proof template as Theorem~4, we have tightened
Theorem~5 by replacing the $\varepsilon n$ Chernoff cutoff with $n^{2/3}$,
which removes $\varepsilon$ from the exponent. The lower bound now sums over
the full range using the series identity with $y = n^{-1/(2n)}$.
The result is tightened to $\Theta(n^{5/2}/\!\log^2 n)$.}

\subsection*{Theorems 6 \& 7 (OneMax and LeadingOnes/LastBest, $M=1$)}

\RC{These proofs can be significantly simplified by observing that before ageing
is triggered, the algorithm behaves exactly as RLS. Let $T'$ be the Opt-IA
runtime; we can write
$\mathbf{E}[T'] \le (\tau + \mathbf{E}[T'])\Pr[A] +
\mathbf{E}[T \mid \lnot A]\Pr[\lnot A]$,
from which $\mathbf{E}[T']$ is at most $(1+o(1))\mathbf{E}[T]$.
Probably it makes sense to unite these two theorems into one with a shorter proof.}

\resp{We have unified Theorems~6 and~7 into a single theorem with a streamlined
proof that follows the reviewer's RLS-reduction argument.
The key observation---that $M=1$ is maintained throughout for both variants on
OneMax and for LastBest on LeadingOnes---is now stated explicitly, and the
decomposition via $\Pr[A]$ is used to derive the $(1+o(1))$ factor.
Theorem~8 (FirstBest on LeadingOnes, $\Theta(n^3)$) is kept separate, as the
RLS reduction does not apply when neutral moves cause $H(x,\mathit{best})$ to
grow.
We also fixed a duplicate label bug between the original Theorems~7 and~8.}


%% ======================================================================
\section*{Updated Figures and Table}
%% ======================================================================

The tightened bounds in Theorems~2, 4, and~5 necessitate an update to
Table~1 (\texttt{table:afl}).  We have also improved the two main figures
in the multimodal-analysis sections.

\subsection*{Table~1: Tightened runtime bounds}

\resp{The two rows of Table~1 corresponding to the exponential-decay
operators have been updated to reflect the new $\Theta$~results:
\begin{itemize}[nosep]
\item FCM with $M_{\text{expoF}}$: \textsc{OneMax} entry changed from
      $O(n^{3/2}\log n),\;\Omega(n^{3/2-2\epsilon})$ to
      $\Theta(n^{3/2}/\!\log n)$; \textsc{LeadingOnes} entry changed from
      $O(n^3/\!\log n),\;\Omega(n^{5/2+\epsilon})$ to
      $\Theta(n^3/\!\log^2 n)$.
\item FCM with $M_{\text{expoHD}}$: \textsc{OneMax} entry changed from
      $O(n^{3/2}\log n),\;\Omega(n^{3/2-2\epsilon})$ to
      $\Theta(n^{3/2}/\!\log n)$; \textsc{LeadingOnes} entry changed from
      $O(n^{5/2+\epsilon}/\!\ln n),\;\Omega(n^{9/4-\epsilon})$ to
      $\Theta(n^{5/2}/\!\log^2 n)$.
\end{itemize}
The $M_{\text{linHD}}$ row and all other rows are unchanged.}

\subsection*{Figure~1: Updated \textsc{Cliff} illustration}

\resp{The \textsc{Cliff}$_d$ figure has been replaced with a clearer version
(\texttt{cliff\_updated.png}) that explicitly labels the local optimum
(40~one-bits, fitness~40) and the global optimum (50~one-bits, fitness~40.5)
for the $n=50$, $d=10$ instance.  The caption has been expanded accordingly.
This addresses the reviewer's general request for improved clarity in the
multimodal-function sections.}

\subsection*{New Figure: \textsc{LeadingTrapJump} landscape diagram}

\resp{We have added a new schematic diagram of the
$\mathrm{LTJ}_{k,\rho}$ landscape (file \texttt{ltj\_diagram.png},
labelled as \texttt{fig:ltj} in the revised manuscript).
The diagram shows the \textsc{LeadingOnes} slope, the junction~$J$, and
the two off-slope structures ($\mathcal{B}_k$ and $\mathcal{D}_\rho$).
A cross-reference to this figure has been added in the prose description
of the function.  This complements the reviewer's comments \textbf{L1--L7}
about the need for clearer notation and structure in the LTJ section.}


%% ======================================================================
\section*{Pages 11--13: Cliff and LeadingTrapJump Proofs}
%% ======================================================================

%% --- CLIFF SECTION (Fixes 1-5) ---

\subsection*{Cliff proof sketch (pp.\,11--12)}

\begin{enumerate}[label=\textbf{C\arabic*.},leftmargin=*]

%% FIX 1
\item \RC{(p.11, l56) $r_0$ and the state description of the form
$(r, \mathrm{HD}(r))$ have not been introduced.}

\resp{We have added an explicit introduction of the indexing convention
$r = r_0, r_0+1, \dots$ for successful fitness improvements on the second
slope, together with the notation $\mathrm{HD}(r)$ for the Hamming distance
after $r$ such improvements. The paragraph preceding
Lemma~3 now defines this notation before using it.}

%% FIX 2
\item \RC{(Lemma~4) Specify that you mean the return in one hypermutation.}

\resp{The statement of Lemma~4 (Exponential Tail for Catch-Up Length) now
explicitly says ``Consider a single application of the hypermutation operator;
let $\xi_z$ be the number of $1\!\to\!0$ flips made \emph{within that one
hypermutation}\,\ldots''}

%% FIX 3
\item \RC{(p.11, r59) I would swap the fraction and $K(\lambda)$ like this:
$\frac{1-\rho}{1-\rho e^{2\lambda}} \eqqcolon K(\lambda)$.}

\resp{Done. All three definition sites of $K(\lambda)$ (Lemma~5, Lemma~7, and
the main proof of Theorem~3) now present the closed-form expression first,
followed by $\eqqcolon K(\lambda)$.}

%% FIX 4
\item \RC{(p.12, l9) Specify that $c$ for $\Theta$ is from the potential.}

\resp{The proof now states explicitly: ``Since the mutation potential satisfies
$M \le c\,n$ (where $c$ is the constant from the linear-HD potential), each
iteration evaluates at most $c\,n = \Theta(n)$ individuals.''}

%% FIX 5
\item \RC{(p.12, l45) Not the buffer itself, but $e^{-\lambda B_t}$.}

\resp{The proof of Corollary~2 (Buffer Barrier Bound) now begins with a
clarification: ``Ville's inequality applies to nonnegative supermartingales,
but the buffer $B_t$ is not one \ldots\ We therefore work instead with the
transformed process $M_t := \alpha^{-t}e^{-\lambda B_t}$, which \emph{is} a
nonnegative supermartingale.'' The remainder of the proof then proceeds via
$M_t$.}

\end{enumerate}

\subsection*{Typographical correction (Cliff, LastBest corollary)}

\resp{While implementing the above revisions we noticed a typo in the
LastBest corollary for \textsc{Cliff}: the denominator read~$d^3$ instead
of~$d^2$.  The prose preceding the corollary correctly derives the
LastBest bound as a factor~$d/n$ better than the FirstBest
bound~$O(n^5 \log(n/d)/d^3)$, giving
$O(n^4 \log(n/d)/d^2 + n^2 \log d)$; the corollary now matches.
We also corrected a related typo in the supplementary material where $n^e$
appeared in place of~$n^4$.}


%% --- LTJ SECTION (Fixes 6-12) ---

\subsection*{LeadingTrapJump proof sketch (p.\,13)}

\begin{enumerate}[label=\textbf{L\arabic*.},leftmargin=*]

%% FIX 6
\item \RC{(p.13, l11) Notation $x_{i>r}$ has not been introduced.}

\resp{We have added the definition
``$x_{i>r} := (x_{r+1}, x_{r+2}, \dots, x_n)$''
immediately after the $\mathrm{LTJ}_{k,\rho}$ function definition.}

%% FIX 7
\item \RC{(p.13, l21) Local optimum $\to$ local optimum $0^r 1^{n-r}$.}

\resp{Both occurrences of ``local optimum'' in the relevant sentence now include
the explicit form $0^{r}1^{(0.9n-r)}$.}

%% FIX 8
\item \RC{(Theorem~11) I do not understand why we should introduce $r$. Is it
not enough to replace it with $k$ everywhere? I do not see the reason for the
prefix of zeros in $\mathcal{B}_k$ to be of odd length.}

\resp{We have added a parenthetical explanation at the point where $r := 2k+1$
is defined: $r$ is odd so that the block-jump prefix
$0^r 1^{(0.9n-r)}\!\ast$ forces an \emph{even} number of bit-flips to reach
$\mathcal{B}_k$, which is the key parameter controlling the per-iteration
success probability $\Theta(n^{-(r+1)}) = \Theta(n^{-(2k+2)})$.
The alias $r$ is introduced as shorthand to avoid writing $2k+1$ in every
exponent throughout the proof.}

%% FIX 9
\item \RC{(p.13, l47) I am not sure I understand what you mean by ``epoch
starting at $J$''. Each epoch starts with a random individual, and only after
that the algorithm goes to $J$.}

\resp{The reviewer is correct. We have replaced all instances of ``Let an ageing
epoch begin at the junction'' with ``Let an ageing epoch \emph{reach} the
junction $J$ (which occurs with high probability by
Lemma~[junction-first-density]; we condition on this event throughout).''
We have also added a clarifying sentence: ``Each ageing epoch begins with a
uniformly random reinitialisation and runs for $\tau$ iterations; by
Lemma~[junction-first-density] it ascends to the junction $J$ with high
probability before reaching any off-slope structure.''}

%% FIX 10
\item \RC{(p.13, r3--r7) But initially there are $L/2$ one-bits, not $L$.}

\resp{Tightened. The proof now states: ``since the $L$-bit prefix of the new
parent is a uniformly random string (from the reinitialisation), it contains
approximately $L/2$ one-bits in expectation; the total number of ones therefore
drops by at most $L/2 \le \tfrac12 \log_2 n$ when these are flipped.''}

%% FIX 11
\item \RC{(p.13, r10) Why $q^\star < 1/4$?}

\resp{We have expanded the justification. The text now shows the explicit
calculation: choosing $\xi$ close to $0$ gives
$q^\star \approx 2c(1-\rho)$; since $\rho > 4/5$ we have $1-\rho < 1/5$, so
$2c(1-\rho) < 2 \cdot \tfrac{5}{8} \cdot \tfrac{1}{5} = \tfrac{1}{4}$ (using
$c < 5/8$). Any fixed $\xi$ small enough to keep $2c(1-\rho+\xi) < 1/4$ then
satisfies all requirements.}

%% FIX 12
\item \RC{(p.13, r40) I am not sure I understood the coupon-collector effect
here and why we need both times.}

\resp{We have expanded the explanation: ``Since $B$ (entering
$\mathcal{B}_k$) and $D$ (entering $\mathcal{D}_\rho$) are mutually exclusive
within a single epoch, the algorithm must discover them in \emph{different}
epochs. By a coupon-collector argument (we need to `collect' both coupons $B$
and $D$, which arrive independently across epochs with probabilities $p_B$ and
$p_D$ respectively), the expected number of epochs until both have been seen is
$O(1/p_B + 1/p_D)$.''}

\end{enumerate}


%% ======================================================================
\section*{Pages 3--10: Local Fixes}
%% ======================================================================

\begin{enumerate}[label=\textbf{P\arabic*.},leftmargin=*]

%% FIX 1
\item \RC{(Table~2 caption, p.\,3) It is not clear what $c$ is (in the
definition of~$\beta$).}

\resp{The caption of Table~2 now reads
``$\beta := \log_2 \tfrac{20}{c}$, where $c\le 1$ is the constant factor in
the $M_{\text{linHD}}$ potential
$M = \lfloor c\cdot H(x,\mathit{best})\rfloor$.''}

%% FIX 2
\item \RC{(p.\,3, r54) Is it assumed that $f_{\mathrm{OPT}}$ is non-negative or
even positive? In this case, due to rounding up, is not it always just~$cn$?}

\resp{The formula description now explicitly states
``$f_{\text{OPT}}\ge 0$ is the optimal fitness value and $v>0$ is the fitness
of the current individual.''  The previous phrasing ``best known fitness''
was misleading since $f_{\text{OPT}}$ denotes the true optimum, not an estimate.}

%% FIX 3
\item \RC{(p.\,3, r56--60) For OneMax a different potential was used: if we
minimize OneMax, then $f_{\mathrm{OPT}}=0$ and the potential is always~$cn$.}

\resp{We have added a clarifying paragraph immediately before the existing
analysis of the original IPH on \textsc{OneMax}: ``Note that when
\textsc{OneMax} is treated as a maximisation problem (as we do throughout this
paper), $f_{\text{OPT}}=n>0$ and the potential decreases as fitness improves.
In contrast, under minimisation ($f_{\text{OPT}}=0$), the formula gives
$M=cn$ always, recovering the static potential---the original motivation for
the modifications we study here.''}

%% FIX 4 + FIX 11 (merged)
\item \RC{(p.\,4, l58) Floor command is missing.
\newline
(p.\,8, r15) Here you use smallcaps in the index (for `LIN'), but everywhere
else you do not do this.}

\resp{Both issues concerned the same display equation defining the practical
$M_{\text{linHD}}$ potential.  We have corrected the subscript from
\texttt{\textbackslash textsc\{linHD\}} to
\texttt{\textbackslash text\{linHD\}} (matching the \texttt{\textbackslash linHD}
macro defined in the preamble), ensured that explicit
$\lfloor\cdot\rfloor$ floor brackets are present, and changed
\texttt{\textbackslash text\{best\}} to
\texttt{\textbackslash mathit\{best\}} for consistency with the algorithm
pseudocode.}

%% FIX 5
\item \RC{(p.\,5, r55) Either add a factor $(1-o(1))$ in front of the sum, or
add a note: ``(w.o.p., thus it is also $\Omega(n^2)$ unconditionally).''}

\resp{We have added the factor $(1-2^{-\Omega(n)})$ in front of the sum in
the lower bound of Theorem~1 and appended the parenthetical note: ``the factor
$(1-2^{-\Omega(n)})$ accounts for the w.o.p.\ Chernoff event; since
$2^{-\Omega(n)}=o(1)$, the bound holds unconditionally.''}

%% FIX 7
\item \RC{(p.\,7, l9) It is not a trivial conclusion about the probability to
skip a level. Could you add more arguments here or add a reference? I assume
it should be in the first paper on LeadingOnes by Rudolph (1997).}

\resp{We have expanded the argument in the proof of Theorem~3.  The text now
reads: ``the probability of not skipping level~$i$ (i.e., the probability that
the number of leading $1$-bits increases by exactly one when the leftmost
$0$-bit is flipped) is $\Omega(1)$.  This holds because the expected number of
free-riders (consecutive $1$-bits following the leftmost $0$-bit) is at most
$\sum_{j\ge1}j\cdot 2^{-(j+1)}<1$, so the level advances by exactly~$1$ with
probability at least $1/2 = \Omega(1)$; see also [DJW02].''
We cite [DJW02] (Droste, Jansen, and Wegener, 2002) where the free-rider
calculation is implicit; this reference is already in our bibliography.}

%% FIX 12
\item \RC{(p.\,9, Algorithms~3 \& 4) I assume there should be a check if the
re-initialized individual is better than the best, even though it is not likely
to happen.}

\resp{Both algorithms now include a best-update check immediately after
reinitialisation.  Algorithm~3 (LastBest) uses
$f(w) \geq f(\text{best})$ and Algorithm~4 (FirstBest) uses
$f(w) > f(\text{best})$, consistent with their respective update rules
($\geq$ for LastBest, $>$ for FirstBest).}

%% FIX 13
\item \RC{(p.\,9, l30) Theorem should work even for
$\tau\ge Cn\ln n$ for some constant~$C$ that is large enough.}

\resp{The rewritten Theorem~6 (which now covers both algorithm variants on
\textsc{OneMax} via the RLS-reduction lemma) already uses the relaxed
threshold $\tau \ge Cn\ln n$ with $C>2$.  We have added a remark after the
proof making the verification explicit: $(1-1/n)^\tau \le n^{-C}$, a union
bound over $n$ levels gives $\Pr[A]\le n^{1-C}=o(1)$ for $C>1$, and
$\tau\cdot\Pr[A]=o(1)$ for $C>2$.}

%% FIX 14
\item \RC{(p.\,10, l38) It should be $\ell<2n/3$ (instead of $\ell\le 2n/3$).}

\resp{Corrected: the proof of Theorem~8 (FirstBest on LeadingOnes) now reads
``Fix a level~$\ell$ with $n/3\le \ell < 2n/3$.''}

%% FIX 15
\item \RC{(p.\,10, l48) It should be $m/(n-1)$ instead of $m/n$.}

\resp{Corrected. The text now reads: ``the expected number of tail-first
iterations is $\frac{m}{n-1}(\mathbb{E}[T]-1) = m = \Theta(n)$ (using that,
conditional on not hitting the improving bit $\ell+1$, the first flip is
uniform over the remaining $n-1$ positions, of which $m$ are tail
positions).''}

%% FIX 16
\item \RC{(p.\,10, l52) Actually, not \emph{at least} once, but
\emph{exactly} once.}

\resp{Corrected: ``a constant fraction of the $m$ tail positions has been
flipped \emph{exactly} once with constant probability (since each
hypermutation samples bits without replacement, each position is visited at
most once per call).''}

%% FIX 17
\item \RC{(p.\,10, l53) `best' is a name of a variable, so there should be no
`the' before it. I actually recommend to rename it to $x_{\mathit{best}}$.}

\resp{We have removed the article ``the'' before \emph{best} and typeset the
variable in italic ($\mathit{best}$) throughout the affected proof.
We opted for the minimal fix rather than a global rename to
$x_{\mathit{best}}$, as the existing notation is already established in
Algorithms~3 and~4 and throughout the other proofs.}

\end{enumerate}

\medskip
\noindent\textbf{Note on fixes already addressed by proof replacements.}\quad
The following reviewer comments from pages~3--10 were automatically resolved
by the complete proof rewrites of Theorems~2, 4, and~5:
\begin{itemize}[nosep]
\item \emph{FIX~6} (p.\,6, $\varepsilon/2$ in Theorem~2 lower bound):
  resolved by the new Theorem~2 proof (theorem2\_replacement.tex), which no
  longer uses the $\varepsilon n$ Chernoff cutoff.
\item \emph{FIX~8} (p.\,7, repeated paragraph in Theorem~4 proof): the
  duplicated computation has been removed in the rewritten proof
  (theorem4\_replacement.tex).
\item \emph{FIX~9} (p.\,7, introduce $\varepsilon$ at start of Theorem~4
  lower bound): the rewritten proof does not use $\varepsilon$ as a free
  parameter, so the issue no longer applies.
\item \emph{FIX~10} (p.\,7, awkward proof opening for Theorem~5): the orphan
  display equation has been replaced by a properly introduced Taylor expansion
  in the rewritten proof (thm5Replacement.tex).
\end{itemize}

\end{document}
